\chapter{Introduction}
\label{chap:introduction}

Sport shooting requires accurate and reliable scoring of targets in order to evaluate athlete
performance during training and competition, as well as to improve the training experience.
Traditionally, paper targets have been scored
manually, a process that is time-consuming, error-prone and unsuitable for real-time feedback.
To overcome these limitations, a variety of automatic scoring systems have been developed over
the years, including solutions based on acoustic triangulation, laser beam interruption and
electronic impact sensing. While these systems offer high accuracy and reliability, they typically
involve specialized hardware and high cost, making them less accessible to amateur
clubs and training environments.

Computer-vision-based scoring systems provide a low-cost alternative by using standard
imaging hardware and image-processing algorithms to detect and score bullet impacts on paper
targets. Advances in image processing and machine learning have increased the feasibility of
such approaches, enabling automatic detection of bullet holes under challenging
conditions. However, the application of computer vision to ISSF-compliant paper targets remains
challenging due to the small dimensions involved, the use of decimal scoring and the high
accuracy requirements imposed by official rules.

This work focuses on the automatic scoring of ISSF 10\,m air rifle and 10\,m air pistol paper
targets using low-cost camera-based systems. The objective is to design a computer-vision-based
scoring system capable of detecting bullet holes, estimating their positions with sufficient
accuracy and computing ISSF-compliant decimal scores in near real time. Particular emphasis is
placed on achieving a balance between robustness to real-world imaging conditions and the
geometric precision required for accurate scoring.

To address this problem, the dissertation explores a hybrid system architecture that combines
learning-based methods and classical image-processing techniques. Convolutional neural networks
are considered for robust target and bullet-hole detection, while geometric refinement methods
are employed to recover accurate center coordinates required for scoring. The system is designed
to operate in controlled indoor shooting ranges using low-cost hardware, without relying on
specialized electronic targets.

The remainder of this document is organized as follows. Chapter~2 analyzes the physical and
regulatory constraints of ISSF paper targets and reviews the state of the art in computer-vision-
based target scoring systems. Chapter~3 presents the proposed system architecture and outlines
the image-processing methods considered for target localization, bullet-hole detection and
geometric refinement. Chapter~4 concludes the report with a summary of the main findings and a
discussion of the expected contributions and future work during the dissertation phase.


